% !TeX root = ../papers/nc-p3.tex
\section{四行列恒等式を持つ全成分の正則性判定}
\label{sec:four-matrix-all-regular}
\begin{lemma}\label{lem:binary-simple-trace-escape}
有限四辺の標準正則窓で、両向きとも頂点との第一Extを持つ均衡二次元単純が存在しないとする。
少なくとも一向きにそのExtを持つ二次元単純は、次数倍率と同型を除いて有限個である。
この補題にはPI性を要しない。
\end{lemma}
\begin{proof}
$S$ から頂点への第一Extがどこかで非零の場合を扱う。仮定より逆向きは全頂点で零である。
薄い単純 $D$ を一つ固定する。補題\ref{lem:periodic-tail-vertex-bridges}の尾部計算は
行列恒等式を使う前に、ある $\lambda\ne0$ で $\operatorname{Ext}^1(D_\lambda,S)\ne0$ を与えた。
従ってこの型の単純の各次数軌道は
\[
   I=\{[S]: S\text{ は単純},\quad
       \text{ある }i\text{ で }\operatorname{Ext}^1(S,k_i)\ne0,\quad
                       \operatorname{Ext}^1(D,S)\ne0\}
\]
に会う。ここで $[S]$ は $2\delta$ のアフィン表現商の指標で、閉軌道の半単純同型型を記録する。
$I$ は構成可能である。単純性は開で、有限両側Koszul複体の階数で各Ext次元を調べ、
有限型の不変式商へ射影するからである。商は正次数の閉路traceの有限個と定数で生成される
（標数零の行列不変式を頂点冪等元とともに適用）。
それらの正次数を $w_j$、座標を $x_j$ とする。倍率はここでは全矢の共通倍率とする。

$I$ が無限なら、その閉包中に一般点が $I$ にあるアフィン曲線を取れる。
その滑らかな射影モデルのある点で少なくとも一つの $x_j$ が極を持つ。
曲線の一般指標を実現する単純 $S$ の行列を有限体拡大後に取る。極の付値を $v$ とし、
必要なら分岐拡大して
\[
             v(u)=\max_{x_j\ne0}(-v(x_j)/w_j)>0
\]
となる元 $u$ を選ぶ。$S$ と $D$ の全矢を同時に $u$ 倍する。
全 $u^{w_j}x_j$ が整で少なくとも一つが単元となる。
この $S$ の整な格子表示が存在する。一般の絶対単純の像は全行列環であるから
経路語の像から体上の基底を取れる。どの二語の積のtraceも整なので、整係数の経路像は
この基底のtrace双対格子に含まれる。よって経路像の環は付値環上有限であり、
一つの自由格子に作用させた和が安定な格子を与える。頂点冪等元で分けて還元し $S_0$ とする。
$D$ には定数基底の格子を使えば、縮小した $D$ の還元は全矢零の頂点単純の和である。

$S_0$ の組成因子には頂点単純がない。全てが頂点なら正次数traceが全て零となり単元の存在に反する。
非薄単純があれば全位置次元が2以上なので、$2\delta$ 全体を占めて $S_0$ 自身が単純となる。
薄いものと頂点単純の混在も、両群間の第一Extが両向きに零なので一般絶対単純に平滑化できない。
実際両群に分裂した特殊表示の非対角項を、完備化で一次コサイクルを厚みごとに消して除ける。
従って特殊単純が薄いものだけなら、それらの全てで頂点からの第一Extは零である。
$S_0$ 自身が二次元単純なら、一般表現から頂点への非零Extは階数の半連続性で特殊化にも残る。
曲線上で頂点は一つに固定できる。補題の仮定から逆向きは全て消える。
いずれにせよ組成列で $\operatorname{Ext}^1(\bigoplus k_i,S_0)=0$ となる。
しかし両表現の同時倍率変更は $\operatorname{Ext}^1(D,S)\ne0$ を保ち、
有限自由なKoszul複体の特殊化でその次元は下がらない。これは矛盾する。
これで $I$ は有限であり、反対向きも反対環で同様である。
\end{proof}

\begin{theorem}\label{thm:four-matrix-all-regular}
任意の成分の有限四辺の標準有限Frobenius候補で、以下は同値である。
\begin{enumerate}
\item 標準Koszul左右AS正則で、各隅は四行四列行列の全恒等式を満たす。
\item 経路長8の四中心元による有限検査が通り、得た有限乗法表で一般の頂点隅の次数が4以下である。
中心剰余は29次、経路次元は30次、左右Koszul複体は32次まででよい。
\end{enumerate}
この枝の全ての正則窓は中心有限な合成整域で、頂点指数は2または4に限られる。
\end{theorem}
\begin{proof}
第一条件を仮定する。もし中心有限でなければ定理\ref{thm:all-three-matrix-clifford}と
定理\ref{thm:four-matrix-balanced-simples}により、全非頂点単純は薄いか、
頂点との第一Extがちょうど片向きにある $2\delta$ である。
薄い単純には有限点グラフにより次数倍率を除き有限個の型しかない。
残る二次元単純も前補題で次数倍率を除いて有限個である。

従って全単純を消すイデアルは有限個の斉次素イデアルの交叉であり、それらの商は全て
GK次元1以下である。実際各非頂点代表の矢を不定元 $v$ 倍すると有限行列環 $M_N(k[v])$ への
斉次表示を得る。その核は倍率軌道全体の共通核で、$k(v)$ 上で絶対単純なので素イデアルである。
次数 $n$ の像は $v^n M_N(k)$ に入る。頂点単純も加えれば全てである。折り畳みはアフィンPIで、
全単純の有限次元性とJacobson根基の冪零性により、これらの素イデアルを並べた積の適当な冪は零である。

このような積はAS正則性に反する。頂点射影に右から順に各イデアルを掛けて、零となる直前に
非零の斉次元 $a$ を選ぶ。$N=aB$ はある一つの軌道商 $R$ 上の加群である。
頂点商なら有限次元の部分加群が生じ、右ASのsocle消滅に反する。
軌道商では正次数の斉次中心元 $z$ が存在し、$R$ は $k[z]$ 上有限である。
これを直接見るには、経路像の有限個の語から全行列環の体上の基底を選び、非恒等な行列中心多項式に代入する。
十分な正長閉路・経路の線型結合で一般行列環全体が張られる（次数倍率を1として、正長イデアルの像は全体）。
代入の全斉次成分も $v$ 変数の行列として中心で、非零な正次数成分は $c v^l 1_N$ となる。
例えば同じ多項式を全ての経路線型結合で代入して $v=1$ の非零値を取ればよい。
$R\subset M_N(k[v])$ はその $k[v^l]$ 部分加群なので有限である。
$N$ も有限で、その $z$-torsionは有限次元次数付き右部分加群ゆえ零である。
従って $0\to N(-l)\xrightarrow{z}N\to N/zN\to0$ で最後は有限次元。
この最後の加群から $B$ への次数付き第0,1 Extの消滅により
$\underline{\operatorname{Hom}}_B(N,B)$ は非零次数移動と同型になる。
$N$ は一生成でこのHomは内部次数で下に有界だから零でなければならず、包含に矛盾する。

よって全ての正則窓は中心有限で合成整域となる。一般の隅は四行列の恒等式を満たし指数は4以下、
定理\ref{thm:index-four-octic-normalization}で八次中心パラメータが存在する。
逆向きと指数検査は同じ定理による。全四元が中心なので
命題\ref{prop:central-arrow-module-cutoff}で $L=4(8-1)=28$ とすれば三上限を得る。
有限辺で指数1はなく指数3も系\ref{cor:no-index-three-finite-noetherian-window}で除かれる。
\end{proof}
